\documentclass[10pt,a4paper]{amsart}
\usepackage[margin=19mm]{geometry}
\usepackage{amsmath,amssymb,mathtools}
\usepackage[scr=rsfs,cal=euler]{mathalfa}
\usepackage{xfrac}
\usepackage[unicode,hidelinks,bookmarks=false]{hyperref}
\newcommand{\ie}{i.e.\ }
\newcommand{\cf}{cf.\ }
\newcommand{\resp}{respectively\ }
\newcommand{\Subsection}{\subsection}
\providecommand{\nobreakdash}{\nobreak\hbox{-}}
\setlength{\emergencystretch}{2em}
\multlinegap0pt
\usepackage{mathtools}
\usepackage{amssymb}
\usepackage[shortlabels]{enumitem}
\newtheorem{theorem}{Theorem}
\newtheorem{lemma}{Lemma}
\newcommand{\N}{\mathbb N}
\DeclareMathOperator{\SG}{SG}
\newcommand{\Sd}{\mathcal S_d}
\newcommand{\Z}{\mathbb Z}
\newcommand{\cL}{\mathcal L}
\newcommand{\cW}{\mathcal W}
\newcommand{\supp}{\operatorname{supp}}
\newcommand{\ulim}[2]{#1\!\operatorname{-lim}_{#2}}
\title{Nil-Bohr Sets and Sums with Bounded Gaps}
\author{Yang Cao, Jianjie Zhao}
\date{Source: arXiv:2608.26568v1 (CC BY 4.0)}
\begin{document}
\maketitle

\noindent\emph{Source and licence.} Nil-Bohr Sets and Sums with Bounded Gaps. Version arXiv:2608.26568v1 (CC BY 4.0), distributed by arXiv under the Creative Commons Attribution 4.0 International licence, http://creativecommons.org/licenses/by/4.0/, as recorded in the arXiv licence metadata for this exact version and shown on the abstract page https://arxiv.org/abs/2608.26568v1. Full licence notice: \texttt{licenses/arxiv-2608.26568-CC-BY-4.0.txt}.\par\smallskip

\noindent\textit{Editorial scope.} Counted unit: the article's theorem thm:dynamical (source0.tex lines 196--205). The article's complete original proof of it is delivered with it (source0.tex lines 1262--1453, one \texttt{proof} environment of 192 lines, which the article places away from the statement). No mathematics is written, repaired or re-derived: the statement and the proof are the article's own lines, in the article's own notation, and no retained line is substituted or re-flowed.  Retained uncounted as the source-local context the proof needs: lines 495--495; lines 788--800; lines 947--947; lines 1090--1093; lines 1181--1197.  Nothing was omitted from any retained span, so the package carries no omission record.  No retained line contains a citation, so the wrapper carries no bibliography and no bibliography entry.  No retained line contains a figure environment, an image-inclusion command or drawing code, so no image is delivered and the package declares no asset dependency.  Editorial formatting only, in addition: an AMS \texttt{amsart} preamble that restates the article's own declarations and macros used by the retained lines, namely the environments \texttt{theorem}, \texttt{lemma} on separate counters, together with the standard packages those lines need.  The retained environments print in this document as Theorem 1, Lemma 1 in order of appearance, whereas the article prints the same items with the numbering of its own full document.  The complete native source of the article as distributed in the arXiv e-print for this exact version is bundled unchanged as \texttt{sources/208641/source0.tex}.
\section{Sums with gaps and paths in the cube}\label{sec:cube}

\begin{lemma}\label{lem:path-realisation}
Let $s=((t_1,\varepsilon_1),\ldots,(t_k,\varepsilon_k))$ be a nonempty admissible timed path from $0_d$ to $0_d$.  Suppose that $x^s(t)\neq0_d$ for
every integer $t$ with $t_1<t<t_k$.  Then there exists a unique
$\alpha\in\Sd$ such that
\[
  \operatorname{Ev}(\alpha)
 =((t_1,\varepsilon_1),\ldots,(t_k,\varepsilon_k)).  
\]
Moreover, 
\[
 T^{p_\alpha}=\prod_{j=1}^k v_{t_j}(\varepsilon_j).   
\]
\end{lemma}

\section{A compact subgroup of the Ellis group}\label{sec:compact-subgroup}

\begin{lemma}\label{lem:word-limit}
Let $w=(\varepsilon_1,\ldots,\varepsilon_k)\in\cW$,  
then $h_{\varepsilon_1}\cdots h_{\varepsilon_k}\in K$.
\end{lemma}

\begin{lemma}\label{lem:cap}
Let $G$ be a group, and let $G_1,\ldots,G_{d+1}$ be subgroups satisfying $[G_j,G_1]\subseteq G_{j+1}$ for every $1\leq j\leq d$.  Assume that these subgroups form the filtration
\[
 G=G_1\supseteq G_2\supseteq\cdots\supseteq G_{d+1}=\{e\}.
\]
Let $C\in G_2$ and $\ell_1,\ldots,\ell_{d-1}\in G_1$.  For
$U=\{i_1<\cdots<i_k\}\subseteq\{1,\ldots,d-1\}$ set
\[
 B_U=\ell_{i_k}\cdots\ell_{i_1},
 \ z_U=B_UCB_U^{-1},
\]
and put $z_\varnothing=C$.  Then
\[
  C\in \big\langle z_U:\varnothing\neq U\subseteq\{1,\ldots,d-1\}\big\rangle.  
\]
 
\end{lemma}

\begin{theorem}\label{thm:dynamical}
Let $(X,T)$ be a compact Hausdorff distal system whose Ellis group admits a
filtration of length $d$.  Then, for every $x\in X$, every
neighbourhood $U$ of $x$, and every sequence $P=(p_i)_{i\geq1}$ of positive
integers, there exists $\alpha\in\Sd$ such that
\[
 T^{p_\alpha}x\in U.
\]
Equivalently, every return time set $N(x,U)$ is an $\SG_d^*$ set.
\end{theorem}

\begin{proof}[Proof of Theorem~\ref{thm:dynamical}]
Assume first that $d\geq2$.  Fix $P=(p_i)_{i\geq1}$, and use the
notation from Sections~\ref{sec:cube} and~\ref{sec:compact-subgroup}.
Let $\pi:E\longrightarrow E/E_2$ be the quotient map.  Inversion is continuous
on the compact topological
group $E/E_2$.  Hence, for every $r\in\Z/d\Z$,
\[
\begin{aligned}
 \pi(L_r)
 &=\ulim{  p_r}{t\in D_r}\pi(Q_t^{-1})\\
 &=\ulim{ p_r}{t\in D_r}\pi(Q_t)^{-1}\\
 &=\left(\ulim{\ p_r}{t\in D_r}\pi(Q_t)\right)^{-1}\\
 &=\pi(R_r)^{-1}.
\end{aligned}
\]
Set $C:=R_0L_0$.
Then $\pi(C)=\pi(R_0)\pi(L_0)=e$,
so $C\in\ker\pi=E_2$. We next show that this element $C$ belongs to $K$.
 

Apply Lemma~\ref{lem:cap} with $G_j=E_j$ and $\ell_i=L_i$.  Let
$U=\{i_1<\cdots<i_k\}$ be a nonempty subset of $\{1,\ldots,d-1\}$, and set $B_U=L_{i_k}\cdots L_{i_1},\ H_U=R_{i_1}\cdots R_{i_k}$.
 
 
Consider
\[
 w_U=(+i_k,\ldots,+i_1,-i_1,\ldots,-i_k)
\]
and
\[
 \widetilde w_U=(+i_k,\ldots,+i_1,-0,+0,-i_1,\ldots,-i_k).
\]
For $A\subseteq\Z/d\Z$, let $\mathbf1_A\in E_d$ be its indicator vector, i.e., 
\[
 (\mathbf1_A)_r= \begin{cases}
                  1, r\in A\\
                  0, r\notin A, 
                  \end{cases}   r\in \Z/d\Z.
\]
The state path of $w_U$ is
\[
 \mathbf1_{\{0\}}
 \xrightarrow{+i_k}\cdots\xrightarrow{+i_1}
 \mathbf1_{\{0\}\cup U}
 \xrightarrow{-i_1}\cdots\xrightarrow{-i_k}
 \mathbf1_{\{0\}}.
\]
The state path of $\widetilde w_U$ is
\[
 \mathbf1_{\{0\}}
 \xrightarrow{+i_k}\cdots\xrightarrow{+i_1}
 \mathbf1_{\{0\}\cup U} \xrightarrow{-0}
 \mathbf1_U\xrightarrow{+0}
 \mathbf1_{\{0\}\cup U} 
 \xrightarrow{-i_1}\cdots\xrightarrow{-i_k}
 \mathbf1_{\{0\}}.
\]
%\[
% \mathbf1_{\{0\}\cup U}\xrightarrow{-0} \mathbf1_U \xrightarrow{+0} \mathbf1_{\{0\}\cup U}.
%\]
Since $U\neq\varnothing$, we have $\mathbf1_U\neq0_d$.  Thus both
paths avoid $0_d$ except at their initial and terminal states, and hence $w_U,\widetilde w_U\in\cW$.
Lemma~\ref{lem:word-limit} implies  
\[
 B_UH_U\in K, \ B_UCH_U\in K.
\]
Since $K$ is a group,
\[
 (B_UCH_U)(B_UH_U)^{-1}=B_UCB_U^{-1}\in K.
\]
This holds for every nonempty $U\subseteq\{1,\ldots,d-1\}$.  
Lemma~\ref{lem:cap} therefore gives $C\in K$.  Since $K$ is a group, we also have $C^{-1}\in K$.
 

Choose $q\in\delta\cL$ with $\Phi(q)=C^{-1}$.
For every $N\in\N$, $\cL_{>N}\in q$.
Let $a,b\in D_0$ and $s\in\cL$ satisfy
\[
 a<\min\supp(s)\leq\max\supp(s)<b.
\]
Write  
\[
 s=((t_1,\varepsilon_1),\ldots,(t_k,\varepsilon_k))
\]
and set
\[
 \widetilde s
 =((a,+0),(t_1,\varepsilon_1),\ldots,(t_k,\varepsilon_k),(b,-0)).
\]
The state path of $\widetilde s$ is
\[
 0_d \xrightarrow{+0} e_0 \xrightarrow{s} e_0\xrightarrow{-0} 0_d.
\]
The middle path corresponding to $s$ is the realization of an admissible word, so it does not visit $0_d$.
By Lemma~\ref{lem:path-realisation}, there exists $ \alpha(a,s,b) \in \Sd$  such that 
\[
 \operatorname{Ev}( \alpha(a,s,b))=\widetilde s,
\]
and
\[
 \begin{aligned}
 T^{p_{ \alpha(a,s,b) }}
 &=v_a(+0)\Phi(s)v_b(-0)\\
 &=Q_a^{-1}\Phi(s)Q_b.
\end{aligned}   
\]
 

For fixed $a$ and $s$, the element $Q_a^{-1}\Phi(s)$ is a power of $T$ and  hence belongs to $\Lambda(E)$.  Since
\[
 \{b\in D_0:b>\max\supp(s)\}\in p_0,
\]
we have
\[
\begin{aligned}
 &\ulim{  p_0}
      {\substack{b\in D_0\\b>\max\supp(s)}}
 Q_a^{-1}\Phi(s)Q_b\\
 & =Q_a^{-1}\Phi(s)
   \left(\ulim{  p_0}{b\in D_0}Q_b\right)\\
 &=Q_a^{-1}\Phi(s)R_0.
\end{aligned}
\]
For fixed $a$,
\[
 \{s\in\cL:\min\supp(s)>a\}\in q.
\]
Left multiplication by $Q_a^{-1}$ and right multiplication by $R_0$ are continuous.  Hence
\[
\begin{aligned}
 &\ulim{q}{s:\,\min\supp(s)>a}
 Q_a^{-1}\Phi(s)R_0\\
 & =Q_a^{-1}
   \left(\ulim{q}{s\in\cL}\Phi(s)\right)R_0\\
 & =Q_a^{-1}C^{-1}R_0.
\end{aligned}
\]
Right multiplication by $C^{-1}R_0$ is continuous, so
\[
 \begin{aligned}
 &\ulim{  p_0}{a\in D_0}
  \ulim{q}{s:\,\min\supp(s)>a}
  \ulim{ p_0}
       {\substack{b\in D_0\\b>\max\supp(s)}}
 Q_a^{-1}\Phi(s)Q_b\\
  = & \left(\ulim{ p_0}{a\in D_0}Q_a^{-1}\right)C^{-1}R_0 
  = L_0C^{-1}R_0 \\
  =&L_0(R_0L_0)^{-1}R_0 =e.
\end{aligned}   
\]
 
Set $\mathcal R=\{T^{p_\alpha}:\alpha\in\Sd\}$, then $e\in\overline{\mathcal R}$.
For $x\in X$ and a neighbourhood $U$ of $x$, set
\[
 \mathcal O_U=\{g\in E:gx\in U\}.
\]
The evaluation map $g\mapsto gx$ is continuous and $e\in\mathcal O_U$, so
$\mathcal O_U$ is a neighbourhood of $e$.  Hence $\mathcal O_U\cap\mathcal R\neq\varnothing$.
Choose $\alpha\in\Sd$ such that
$T^{p_\alpha}\in\mathcal O_U$.  Then $T^{p_\alpha}x\in U$,
and hence $p_\alpha\in N(x,U)$.
 
 
It remains to treat the case $d=1$, where no nontrivial commutator
cancellation is needed.
Now let $d=1$.  Then $E_2=\{e\}$, so $E$ is a compact abelian topological
group.  Continuity of inversion implies
\[
 L_0=\ulim{  p_0}{a\in D_0}Q_a^{-1}
 =\left(\ulim{  p_0}{a\in D_0}Q_a\right)^{-1}=R_0^{-1}.
\]
For $a<b$ in $D_0=\N$, set
\[
 \alpha=[a,b-1]\cap\N.
\]
Then $\alpha\in\mathcal S_1$ and
\[
 T^{p_\alpha}=Q_a^{-1}Q_b.
\]
Since $\{b\in D_0:b>a\}\in p_0$,
\[
 \ulim{ p_0}{a\in D_0}
 \ulim{  p_0}{\substack{b\in D_0\\b>a}}Q_a^{-1}Q_b
  =L_0R_0 =e.
\]
Therefore
\[
 e\in
 \overline{\{T^{p_\alpha}:\alpha\in\mathcal S_1\}}.
\]
 The same evaluation argument as above now completes the proof.
\end{proof}

\end{document}
